\section*{Personal y Organización de la CNEN}

La Comisión Nacional de Seguridad Nuclear y Salvaguardias es la autoridad que reglamenta las actividades desarrolladas para la utilización de la energía nuclear con fines pacíficos.

El personal de la Comisión está organizado en:

\begin{itemize}
    \item Gerencia de Seguridad Nuclear.
    \begin{itemize}
        \item Departamento de Evaluación.
        \item Departamento de Verificación operativa.
        \item Departamento de Seguridad Física y Salvaguardias.
    \end{itemize}
    \item Gerencia de Tecnología, Reglamentación y Servicios.
    \begin{itemize}
        \item Departamento de Reglamentos, Documentación y Capacitación.
        \item Departamento de Tecnología e Informática.
    \end{itemize}
    \item Gerencia de Seguridad Radiológica.
    \begin{itemize}
        \item Departamento de Supervisión operativa.
        \item Departamento de Vigilancia Radiológica Ambiental.
    \end{itemize}
    \item Comité de Emergencia.
    \item Consultoría Jurídica.
    \item Oficina de Asuntos Internacionales.
    \item Unidad de Finanzas y Administración.
\end{itemize}

\chapter{Análisis de Incidentes y Accidentes}

Estos se clasifican en:

\begin{itemize}
    \item Falla humana.
    \item Falla o carencia del equipo de protección radiológica.
    \item Improvisaciones.
\end{itemize}

Una vez que el problema se presenta y la fuente está fuera de control, las acciones que se sigan deberán programarse de tal manera que el personal involucrado reciba la mínima dosis posible. Debemos recordar que en la mayoría de los accidentes de radiografía industrial (o en la ocurrencia de la perdida del contenedor de la fuente) no hay por medio vidas que pudieran estar en peligro, por lo que hay tiempo suficiente para programar las actividades que deben realizarse. Actuar precipitadamente por lo general da malos resultados.

A continuación se numeran los casos más frecuentes y se esquematiza en forma breve el procedimiento a seguir:

\begin{itemize}
    \item Caída del contenedor.
    \item Sobrexposiciones inadvertidas.
    \item Fuente suelta.
    \item Accidentes automovilísticos.
    \item Pérdida o robo del contenedor.
\end{itemize}
